\documentclass[11pt,a4paper]{article}
\usepackage[utf8]{inputenc}
\usepackage[T1]{fontenc}
\usepackage[spanish,es-noshorthands]{babel}
\usepackage[margin=2.5cm]{geometry}
\usepackage{amsmath,amssymb,mathtools}
\usepackage{enumitem}
\usepackage{booktabs}
\usepackage{tikz}
\usepackage{pgfplots}
\pgfplotsset{compat=1.18}

\newcounter{ejnum}
\newenvironment{ejercicio}{\refstepcounter{ejnum}\par\medskip\noindent\textbf{Ejercicio \theejnum.}~}{\par}
\newenvironment{solucion}{\par\noindent\textit{Solución.}~}{\par\vspace{1em}}

\title{Práctica 3: Ley de los grandes números y teorema central del límite \\ \large Unidad 4: Función generadora de momentos}
\author{Probabilidad y Estadística (5765) \\ Universidad Nacional del Sur}
\date{}

\begin{document}
\maketitle

\begin{ejercicio}
El tiempo de vida útil de una lamparita (en horas) tiene media $\mu=1000$ y desviación estándar $\sigma=200$ (distribución desconocida). Se toma una muestra de $n=64$ lamparitas independientes y se calcula el promedio muestral $\overline X_{64}$. Usando Chebyshev, acotar $P(|\overline X_{64}-1000|\geq 50)$.
\end{ejercicio}

\begin{solucion}
Sabemos $E(\overline X_{64})=\mu=1000$ y $\text{Var}(\overline X_{64}) = \sigma^2/n = 200^2/64 = 40000/64 = 625$.

Por Chebyshev con $k=50$:
\[
P(|\overline X_{64}-1000|\geq 50) \leq \frac{625}{50^2} = \frac{625}{2500} = 0.25.
\]
La probabilidad de que el promedio muestral se aparte más de $50$ horas de la media real es a lo sumo $25\%$.
\end{solucion}

\begin{ejercicio}
Retomando el ejercicio anterior: ¿cuántas lamparitas $n$ habría que muestrear para garantizar, mediante Chebyshev, que $P(|\overline X_n - 1000|\geq 50) \leq 0.05$?
\end{ejercicio}

\begin{solucion}
Por Chebyshev, $P(|\overline X_n-1000|\geq 50) \leq \dfrac{\sigma^2}{n\cdot 50^2} = \dfrac{40000}{2500\,n} = \dfrac{16}{n}$.

Queremos $\dfrac{16}{n}\leq 0.05 \iff n \geq \dfrac{16}{0.05} = 320$.

\textbf{Se necesitan al menos $n=320$ lamparitas.} Esta es una cota conservadora (por la generalidad de Chebyshev); usando el TCL se puede refinar la estimación si se conoce mejor la forma de la distribución.
\end{solucion}

\begin{ejercicio}
Una máquina fabrica tornillos, y la probabilidad de que un tornillo sea defectuoso es $p=0.04$. Se produce un lote de $n=500$ tornillos independientes. Sea $\hat p$ la proporción muestral de defectuosos.
\begin{enumerate}[label=(\alph*)]
\item Calcular $E(\hat p)$ y $\text{Var}(\hat p)$.
\item Usando el TCL, aproximar $P(0.03 \leq \hat p \leq 0.06)$.
\end{enumerate}
\end{ejercicio}

\begin{solucion}
\textbf{(a)} Como $\hat p = \overline X_n$ con $X_i\sim\text{Bernoulli}(0.04)$ i.i.d.:
\[
E(\hat p) = p = 0.04, \qquad \text{Var}(\hat p) = \frac{p(1-p)}{n} = \frac{0.04(0.96)}{500} = \frac{0.0384}{500} = 0.0000768.
\]
$\sigma_{\hat p} = \sqrt{0.0000768} \approx 0.008764$.

\medskip
\textbf{(b)} Por el TCL, $\hat p \ \dot\sim\ N(0.04,\ 0.0000768)$. Estandarizando:
\[
P(0.03\leq\hat p\leq 0.06) \approx \Phi\left(\frac{0.06-0.04}{0.008764}\right) - \Phi\left(\frac{0.03-0.04}{0.008764}\right) = \Phi(2.282) - \Phi(-1.141).
\]
\[
\approx 0.9888 - 0.1269 = 0.8619.
\]
La probabilidad aproximada de que la proporción muestral de defectuosos esté entre $3\%$ y $6\%$ es $86.2\%$.
\end{solucion}

\begin{ejercicio}
Un ascensor tiene una capacidad máxima de carga de $2000$ kg. El peso de las personas que lo usan tiene media $\mu=75$ kg y desviación estándar $\sigma=12$ kg. Si suben $n=25$ personas independientes, aproximar la probabilidad de que se supere la capacidad máxima.
\end{ejercicio}

\begin{solucion}
Sea $S_{25} = \sum_{i=1}^{25} X_i$ el peso total. Por el TCL, $S_{25} \ \dot\sim\ N(n\mu, n\sigma^2) = N(25\cdot75,\ 25\cdot144) = N(1875, 3600)$, con $\sigma_{S}=\sqrt{3600}=60$.

\medskip
\[
P(S_{25} > 2000) \approx P\left(Z > \frac{2000-1875}{60}\right) = P(Z>2.083) \approx 1-\Phi(2.08) \approx 1-0.9812 = 0.0188.
\]
Hay aproximadamente un $1.9\%$ de probabilidad de que el peso total supere los $2000$ kg con $25$ personas.
\end{solucion}

\begin{ejercicio}
Una empresa de telefonía estima que el $65\%$ de sus clientes renueva el plan anual. Se consulta a una muestra aleatoria de $n=150$ clientes independientes. Sea $X$ el número de clientes que renueva.
\begin{enumerate}[label=(\alph*)]
\item Verificar que se cumplen las condiciones para aproximar $X$ por una normal.
\item Aproximar, con corrección por continuidad, $P(X \geq 100)$.
\end{enumerate}
\end{ejercicio}

\begin{solucion}
\textbf{(a)} $X\sim\text{Bin}(150,0.65)$. Verificamos: $np = 150(0.65)=97.5 \geq 5$ y $n(1-p)=150(0.35)=52.5\geq5$. Se cumplen ampliamente las condiciones para la aproximación normal.

\medskip
\textbf{(b)} $\mu = np = 97.5$, $\sigma^2 = np(1-p) = 150(0.65)(0.35) = 34.125$, $\sigma \approx 5.842$.

Con corrección por continuidad, $P(X\geq 100) = P(X\geq 99.5)$ en la escala continua:
\[
P(X\geq100) \approx P\left(Z \geq \frac{99.5-97.5}{5.842}\right) = P(Z\geq 0.3424) \approx 1-\Phi(0.34) \approx 1-0.6331 = 0.3669.
\]
La probabilidad aproximada de que al menos $100$ de los $150$ clientes renueven es $36.7\%$.
\end{solucion}

\begin{ejercicio}
El número de accidentes de tránsito diarios en una avenida sigue una distribución de Poisson con $\lambda=36$ (es decir, en promedio $36$ accidentes por día, sumando todos los cruces de una ciudad grande). Aproximar, usando la normal con corrección por continuidad, la probabilidad de que en un día dado ocurran entre $30$ y $40$ accidentes (inclusive).
\end{ejercicio}

\begin{solucion}
Como $\lambda=36$ es razonablemente grande, aproximamos $X\ \dot\sim\ N(36,36)$, con $\sigma=6$.

Con corrección por continuidad, el evento $30\leq X\leq 40$ se aproxima por $29.5 \leq Y \leq 40.5$:
\[
P(30\leq X\leq 40) \approx \Phi\left(\frac{40.5-36}{6}\right) - \Phi\left(\frac{29.5-36}{6}\right) = \Phi(0.75)-\Phi(-1.0833)
\]
\[
\approx 0.7734 - 0.1393 = 0.6341.
\]
La probabilidad aproximada es de $63.4\%$.
\end{solucion}

\begin{ejercicio}
Se lanza una moneda equilibrada $n=400$ veces. Sea $X$ el número de caras.
\begin{enumerate}[label=(\alph*)]
\item Calcular $E(X)$ y $\text{Var}(X)$.
\item Aproximar $P(190 \leq X \leq 210)$ usando el TCL con corrección por continuidad.
\item Comparar con la cota que se obtendría usando únicamente la desigualdad de Chebyshev.
\end{enumerate}
\end{ejercicio}

\begin{solucion}
\textbf{(a)} $X\sim\text{Bin}(400,0.5)$: $E(X)=np=200$, $\text{Var}(X)=np(1-p)=400(0.25)=100$, $\sigma=10$.

\medskip
\textbf{(b)} Con corrección por continuidad, $P(190\leq X\leq 210)\approx P(189.5\leq Y\leq 210.5)$:
\[
\Phi\left(\frac{210.5-200}{10}\right)-\Phi\left(\frac{189.5-200}{10}\right) = \Phi(1.05)-\Phi(-1.05) = 2\Phi(1.05)-1
\]
\[
\approx 2(0.8531)-1 = 0.7062.
\]
La probabilidad aproximada (TCL) es $70.6\%$.

\medskip
\textbf{(c)} Por Chebyshev, el evento $|X-200|<10.5$ (aproximadamente $190\leq X\leq210$, sin corrección) da, con $k=10.5$:
\[
P(|X-200|<10.5) \geq 1-\frac{100}{10.5^2} = 1-\frac{100}{110.25} \approx 1-0.9070=0.0930.
\]
Chebyshev solo garantiza que la probabilidad es \textbf{al menos} $9.3\%$, una cota mucho más floja e imprecisa que la aproximación normal ($70.6\%$). Esto ilustra por qué, cuando se dispone de suficiente información sobre la distribución (como en este caso, donde $X$ es binomial y se puede aplicar el TCL), conviene usar la aproximación normal en lugar de Chebyshev: esta última es válida siempre, pero da cotas muy conservadoras.
\end{solucion}

\begin{ejercicio}
Un estudio afirma que el tiempo que un usuario pasa en una aplicación tiene media $\mu=15$ minutos y desviación estándar $\sigma=5$ minutos. Se registra el tiempo de $n=100$ usuarios independientes. Aproximar $P(\overline X_{100} > 16)$, e interpretar el resultado en el contexto de la ley de los grandes números.
\end{ejercicio}

\begin{solucion}
Por el TCL, $\overline X_{100} \ \dot\sim\ N\left(15,\ \dfrac{25}{100}\right) = N(15, 0.25)$, con $\sigma_{\overline X}=0.5$.

\[
P(\overline X_{100}>16) \approx P\left(Z>\frac{16-15}{0.5}\right) = P(Z>2) \approx 1-\Phi(2) \approx 1-0.9772=0.0228.
\]
La probabilidad de que el promedio muestral supere los $16$ minutos es aproximadamente $2.3\%$.

\medskip
\textbf{Interpretación:} este resultado es consistente con la ley de los grandes números (Clase 17): al aumentar $n$, la distribución de $\overline X_n$ se concentra cada vez más alrededor de $\mu=15$ (su varianza $\sigma^2/n$ decrece), por lo que valores de $\overline X_n$ apreciablemente distintos de $15$ (como superar $16$) se vuelven cada vez menos probables. El TCL agrega, además, la información de \emph{cómo} se distribuyen esas fluctuaciones (aproximadamente normales), permitiendo calcular la probabilidad exacta de la desviación en lugar de solo acotarla.
\end{solucion}

\end{document}
